\subsection{Vanishing components for balls and specified surfaces}

\begin{theorem}\label{thm:low-concordance-components}
With the bottom and side boundary fixed and the upper face free,
\[
                         \pi_0P_{\mathrm{PL}}(B)=0
\]
in each of the following cases: a compact manifold of dimension zero
or one; a PL ball of any dimension; a compact connected orientable surface other
than $S^2$; or a closed connected nonorientable surface other than
$\mathbb{RP}^2$.  The assertion also holds for finite disjoint unions
of these manifolds and for their stipulated definable models over every
real closed field.
This is a statement about components.  The surface clauses do not
include $S^2$, $\mathbb{RP}^2$, or nonorientable surfaces with boundary.
\end{theorem}

\begin{proof}
We will combine PL isotopies by means of the piecewise linear function
\[
              r(s,t)=\max(0,s+t-1)\qquad(s,t\in I).
\]
It vanishes on the bottom and left edges, and restricts to $t$ and $s$
on the right and top edges.  Composition with this function keeps total
parameter maps PL; the product $st$ would not have that property.

\emph{Balls.}
Let $h\in P_{\mathrm{PL}}(D^m)$, and denote its upper-face restriction
by $a$.  This homeomorphism fixes $\partial D^m$.  There is a jointly
PL Alexander isotopy $a_u$ from the identity to $a^{-1}$ relative to
that boundary \cite[Theorem~5.2, pp.~95--96]{Epstein1966}.
One can see joint PL regularity directly on $D^m=[-1,1]^m$:
\[
 a_u(x)=
 \begin{cases}
  u\,a^{-1}(x/u),&\lVert x\rVert_\infty\leq u,\quad u>0,\\
  x,&\lVert x\rVert_\infty\geq u,
 \end{cases}
 \qquad a_0(x)=x.
\]
On a cell where $a^{-1}(y)=Ay+b$, the first expression is $Ax+ub$.
Its domain is the cone over that cell at height one.  These finitely
many cones and the outer region give a finite PL map; the formulas
match because $a$ fixes the boundary.  The analogous construction with
$a$ gives the inverse family.

The family
\[
              A_s(x,t)=(a_{r(s,t)}(x),t)
\]
starts at the identity and consists of concordances.  Its endpoint
makes $A_1h$ fix the whole boundary of the PL ball $D^m\times I$.
A second Alexander isotopy, now relative to that whole boundary,
connects $A_1h$ to the identity.  Concatenation proves the assertion for
balls, including dimension zero.

\emph{The circle.}
The only remaining connected one-manifold is $S^1$.  A concordance
of the annulus $S^1\times I$ has an orientation-preserving upper
restriction.  This circle homeomorphism admits a jointly PL isotopy
to the identity.  For an explicit construction, lift it to an
increasing PL map $a:\mathbb R\to\mathbb R$ with
$a(x+1)=a(x)+1$.  Subdivide one period at its breakpoints, triangulate
each resulting rectangle in the parameter strip, and extend the
bottom identity and the top map affinely on those triangles.  On each
horizontal piece, the slope is either $1$ or the positive slope of
$a$ on the corresponding interval.  Each horizontal slice is therefore
increasing, and periodicity makes the construction descend to the
circle.  Using $r(s,t)$ as above adjusts the upper restriction without
moving the bottom.

We may consequently assume that a given annulus concordance $h$ fixes
both boundary circles.  Let $\alpha=\{0\}\times I$.  Its image has an
integer winding $k$ relative to the endpoints, measured in
$\mathbb R\times I$.  The homeomorphism
\[
                   T_{-k}([x],t)=([x-kt],t)
\]
removes that winding and itself lies in the identity component of the
concordance group, by
\[
                   T_{-k,s}([x],t)=([x-kr(s,t)],t).
\]
The upper circle is allowed to move during this isotopy.  The adjusted
arc is homotopic to $\alpha$ relative to its endpoints, so Epstein's
relative arc theorem supplies a PL ambient isotopy fixed on both
boundary circles carrying it to $\alpha$
\cite[Theorem~3.1 and proof, pp.~89--90]{Epstein1966}.
Correcting the arc image does not yet fix its parameterization.
In a rectangular neighborhood $(t,z)\in I\times[-1,1]$ of the arc,
let $\lambda$ be the remaining interval homeomorphism and choose
a jointly PL interval isotopy $c_s$ from the identity to
$\lambda^{-1}$, fixed at both endpoints. The family
\[
 \Phi_s(t,z)=\bigl(c_{\max(0,s-|z|)}(t),z\bigr)
\]
is jointly PL, fixes the rectangle boundary, and extends by the
identity outside the rectangle. Its inverse uses the inverse
interval family. At its endpoint it corrects the central arc
pointwise; this is the regular-neighborhood step in the argument
on p.~91 of the same source.

The annular map is orientation preserving, so after this correction
it preserves the two sides of the arc. Cutting along the arc gives
a disk homeomorphism fixed on each copy of the cut arc and on the
remaining boundary. Its relative Alexander isotopy descends when
the two cut edges are glued. Thus a winding number does not survive
when the upper face is free.

\emph{Orientable surfaces.}
Let $F$ be compact, connected, orientable, and different from $S^2$.
Waldhausen works in the PL category and imposes exactly these surface
hypotheses in his product-line-bundle section
\cite[pp.~56 and~65]{Waldhausen1968}.
His Lemma~3.5, including its proof on p.~69, changes a
bottom-and-side-fixed homeomorphism of $F\times I$ by an isotopy
constant on the full boundary into a level-preserving map
\[
                         (x,t)\longmapsto(h_t(x),t).
\]
Here $h_0=\operatorname{id}$ and every $h_t$ fixes $\partial F$.
The jointly PL family
\[
                         (x,t,s)\longmapsto(h_{r(s,t)}(x),t,s)
\]
joins the identity to this level-preserving concordance.  Its inverse
uses the inverse level-preserving family.  This proves the orientable
surface clause.

\emph{Closed nonorientable surfaces.}
Now let $F$ be closed, nonorientable, and different from
$\mathbb{RP}^2$, and put $V=F\times I$.
We verify the hypotheses of the nonorientable three-manifold theorem.
The universal cover of $F$ is the plane and $\pi_1F$ is torsion free
\cite[Corollary~1.8, pp.~85--86; Lemma~4.3, p.~92]{Epstein1966}.
Thus the universal cover of $V$ is $\mathbb R^2\times I$.
A PL sphere in the interior lifts to a sphere bounding a PL ball in
this slab.  Distinct translates of its lifted sphere are disjoint,
so their bounded balls are either disjoint or nested.  Proper
discontinuity and infinite order of a nonidentity deck transformation
exclude nesting: after possibly taking its inverse, its iterates would
put an infinite orbit in one compact ball.  The ball therefore projects
injectively to $V$.  This proves irreducibility.

A two-sided embedded $\mathbb{RP}^2$ in $V$ would lift to the
orientable slab, since its fundamental group maps trivially to the
torsion-free group $\pi_1V$.  A two-sided surface in an orientable
three-manifold is orientable, a contradiction.
The middle surface $F\times\{1/2\}$ is two-sided and
$\pi_1$-injective.  It cannot bound a contractible submanifold because
its nontrivial fundamental group injects into that of $V$.
Consequently $V$ is $\mathbb{RP}^2$-irreducible and sufficiently large
in the precise sense of
\cite[Chapter~II, Definitions~1.2--1.5, pp.~15--17]{Laudenbach1974}.

Write a concordance as $h(x,t)=(u(x,t),v(x,t))$.  It is properly
homotopic to the identity in the boundary-pair sense: first change the
second coordinate by
\[
             (x,t)\longmapsto
             \bigl(u(x,t),(1-s)v(x,t)+st\bigr),
\]
and then change the first coordinate by
\[
             (x,t)\longmapsto\bigl(u(x,(1-s)t),t\bigr).
\]
Throughout these homotopies each boundary component maps to itself and
interior points map to the interior.  These are homotopies of maps;
there is no claim that their intermediate maps are homeomorphisms.

Laudenbach's Theorem~II.6.1 and proof, pp.~46--47, turn this proper
homotopy class into an isotopy class.  To use its smooth formulation
for the given PL endpoints, use the category comparisons in
Chapter~I, Section~1, pp.~7--9: a PL homeomorphism is TOP isotopic to
a smooth one, and the map from PL to TOP isotopy classes is
injective.  Composing the TOP isotopies and then applying this
injectivity gives a jointly PL isotopy $H_s$ from $h$ to the identity.
This argument does not require an auxiliary smooth map to be PL in a
fixed triangulation.

Finally enforce the fixed bottom.  Each $H_s$ preserves the two boundary
components individually.  Write $b_s$ for its restriction to the
bottom, identified with $F$, and replace it by
\[
                   \widehat H_s=(b_s^{-1}\times\operatorname{id}_I)H_s.
\]
The inverse of the total PL boundary family is PL, so this remains a
jointly PL isotopy.  It fixes the bottom, and its endpoints remain $h$
and the identity because $b_0=b_1=\operatorname{id}$.

The fixed bottom also prevents permutations of connected components.
Finite disjoint unions therefore give products.  The relative
PL--definable component comparison transfers the result, with the side
boundary included in the fixed subpolyhedron, to every real closed
field and the stipulated definable category.
\end{proof}


\begin{corollary}\label{cor:all-closed-surface-components}
For every closed PL surface, the bottom-fixed, top-free concordance
component group is zero. This includes disconnected surfaces and
their stipulated definable models over every real closed field.
\end{corollary}

\begin{proof}
Every connected closed surface is orientable or nonorientable.
The preceding theorem covers all such surfaces except $S^2$ and
$\mathbb{RP}^2$. Theorems~\ref{thm:standard-sphere-components}
and~\ref{thm:projective-plane-components} give those two cases.
The fixed bottom makes a finite disjoint union a product on both
groups and component relations.
\end{proof}
